\section{Related Work}
\label{sec:related}

\paragraph{Agent memory and long-horizon retrieval.}
Generative Agents, MemoryBank, and MemGPT organize observations and experience through reflection, updating, or tiered context management~\citep{park2023generative,zhong2024memorybank,packer2023memgpt}.
Mem0, A-MEM, and HippoRAG develop extracted memories, linked records, and graph-based retrieval~\citep{chhikara2025mem0,xu2025amem,gutierrez2024hipporag}.
Graphiti maintains event-derived relations with temporal validity and provenance for retrieval~\citep{zep2026graphiti}.
\system{} derives recurring conditional patterns from daily trajectories, assesses their historical support, and matches applicable patterns to incoming events.

\paragraph{Proactive decision and memory.}
ProactiveAgent studies whether an agent should propose a task or remain silent and uses annotated event data to elicit this capability~\citep{lu2025proactiveagent}.
ContextAgent combines sensory contexts with historical personas to anticipate service needs~\citep{yang2025contextagent}, while PASK couples streaming demand detection with hybrid long-term memory~\citep{xie2026pask}.
Memory-oriented proactivity also supports internal reasoning and execution: MemCog integrates proactive memory exploration into conversational reasoning~\citep{li2026memcog}, and Remember When It Matters uses a separate memory agent to inject selected reminders into long-horizon action agents~\citep{wu2026remember}.
\system{} focuses on constructing service-relevant temporal memory through source checking, historical replay, and current-event matching.

\paragraph{Evaluation of proactive service.}
LoCoMo and LongMemEval assess long-term conversational memory and reasoning~\citep{maharana2024locomo,wu2025longmemeval}.
Proactive benchmarks examine task proposals, latent-need detection, or memory triggering~\citep{lu2025proactiveagent,xie2026pask,li2026memcog}.
ProactiveStream evaluates sequential \textsc{Intervene}/\textsc{Abstain} decisions as agents accumulate memory within isolated, causally replayed event streams, using episode-aware references to assess intervention timing.
Research on notification timing further shows that interruption cost depends on task phase and receptivity~\citep{adamczyk2004ifnotnow,iqbal2007breakpoints,iqbal2008notification,mehrotra2016myphone}.
This motivates evaluating opportunity recovery together with unnecessary intervention.
